\section{Use Cases Revisited}\label{sec:use-cases-revisited}

Section~\ref{sec:use-cases} named four workflows; each is addressed by the delivered system, within the scope this
thesis claims.

\topic{Offline batch composition.}
The compiler emits byte-reproducible Format~1 MIDI with no runtime at playback.
Motivo Studio adds authenticated storage and download, so source can be edited, compiled, and imported into any standard
host.

\topic{Compact symbolic descriptions.}
The synthetic benchmark fixtures expand a handful of source lines into up to 20,000 note events in interactive compile
times on a local Release build.
The drum-pattern example in Chapter~\ref{ch:evaluation} shows the same pattern at smaller scale: rules replace
hand-drawn events, and lowering materializes the timeline.

\topic{Teaching algorithmic composition.}
Motivo Studio pairs diagnostics, piano-roll preview, and browser playback without a local synthesis install.
The integration is tested automatically; classroom or usability studies are outside this thesis.

\topic{Rapid prototyping before DAW polish.}
Patterns and control flow support motif sketches in source, with MIDI export for mixing and micro-editing in a DAW.
This matches Section~\ref{sec:motivation}: Motivo front-loads structure rather than replacing production tooling.

Taken together, these cases show the delivered system meets the offline, symbolic, complementary scope set out in
Chapter~\ref{ch:introduction}.
